\documentclass[acmlarge,11pt,screen,nonacm,natbib=false]{acmart}
\geometry{letterpaper,margin=1in,headheight=14pt,headsep=18pt,footskip=24pt}
\settopmatter{printacmref=false,printccs=false,printfolios=true}
\setcopyright{none}
\usepackage{xurl}
\AtBeginEnvironment{thebibliography}{\raggedright}
\acmDOI{}
\acmISBN{}
\acmYear{}
\renewcommand{\bibliofont}{\normalfont\normalsize}
\newcommand{\own}{\operatorname{own}}
\newcommand{\Succ}{\operatorname{Succ}}
\newcommand{\last}{\operatorname{last}}
\newcommand{\up}{\operatorname{up}}
\newcommand{\nxt}{\operatorname{next}}
\newcommand{\pay}{\operatorname{pay}}
\newcommand{\Hist}{\mathcal H}
\newtheorem{theorem}{Theorem}
\newtheorem{lemma}[theorem]{Lemma}
\newtheorem{corollary}[theorem]{Corollary}
\theoremstyle{remark}
\newtheorem{remark}[theorem]{Remark}
\title[Arena-Independent Memory for Reachability SPEs]{Arena-Independent Memory for Pure Subgame-Perfect Equilibria in Reachability Games}
\author{Zijian Zeng}
\affiliation{\institution{UCSI University}\country{Malaysia}}
\email{1002266693@ucsiuniversity.edu.my}
\renewcommand{\shortauthors}{Zijian Zeng}
\hypersetup{pdfauthor={Zijian Zeng},pdftitle={Arena-Independent Memory for Pure Subgame-Perfect Equilibria in Reachability Games}}
\begin{document}
\begin{abstract}
\normalsize
In a turn-based game with $n$ players and Boolean reachability objectives,
every pure subgame-perfect equilibrium can be replaced by one with exactly
the same winner set using at most $n3^n$ memory states per player. The bound
holds on every set-sized directed arena without dead ends, with arbitrary
branching, and uses vertex updates without storing a previous vertex.
Subgame perfection holds after every legal history, including histories
inconsistent with the profile, with earlier target visits retained.
This answers the arena-independent constrained-memory question posed by Main.

The construction associates finitely many credible payoff labels with each
vertex and visited-target set. It retains a compatible payoff promise and
otherwise resets it by minimizing the previous owner's payoff bit. A
vertex-dependent progress selector fulfills a retained promise without
storing a path position or rank. A losing-bit invariant rules out arbitrary
unilateral deviations. Combining this compression with the existing
lower-semicontinuous-payoff existence theorem yields unconditional pure
finite-memory SPE existence on the same arena class. A finite degree-two
family with $n=k+2$ players forces at least
$\binom{k}{\lfloor k/2\rfloor}\geq 2^k/(k+1)$ states in every SPE satisfying
a single required winner. The bounds determine exponential-order memory
dependence on the player count, without determining the optimal base.
\end{abstract}
\maketitle
\hypersetup{pdfauthor={Zijian Zeng}}
\thispagestyle{plain}
\clearpage

\section{The memory question and the contribution}

In a multiplayer graph game, a player chooses the next vertex at each vertex
it controls. Its reachability objective asks that a target set be visited
at least once. A Nash equilibrium excludes profitable unilateral deviations
from the initial history. A subgame-perfect equilibrium (SPE) must do so
after every legal history, including histories its prescribed strategies
never produce. Thus a succinct equilibrium outcome is insufficient:
the strategies also need credible continuations after deviations and after
different histories merge at the same vertex.

Main's Automata Exchange question asks whether an SPE satisfying a required
subset of reachability objectives, whenever one exists, can be implemented
with memory bounded independently of the arena size~\cite{MainQuestion25}.
We prove a uniform bound, preserving the entire payoff rather than only the
required winners. No finiteness or branching assumption is imposed on the arena.

\begin{theorem}[Payoff-preserving compression]\label{thm:main}
Let $n\geq1$. In every set-sized $n$-player turn-based directed arena without
dead ends, with an initial vertex and Boolean reachability targets, every
pure SPE has a pure SPE with exactly the same winner set and at most $n3^n$
vertex-update memory states per player. The players may share their memory
set, initial state and update function.
\end{theorem}

If an SPE wins for every player in a specified set $S$, Theorem~\ref{thm:main}
gives such an SPE with the stated bound, for every $S$, allowing additional
winners. In particular it resolves the full constrained-memory question.
The theorem separates internal memory from the size of the graph even when
the latter is infinite. Its vertex-dependent strategy tables need not have
a finite effective description on an infinite arena.

\begin{corollary}[Unconditional existence]\label{cor:existence}
Every game in Theorem~\ref{thm:main} admits a pure SPE with at most $n3^n$
vertex-update states per player.
\end{corollary}

This corollary combines the compression with an old unrestricted existence
theorem of Flesch, Kuipers, Mashiah-Yaakovi, Schoenmakers, Solan and
Vrieze~\cite{FKMSSV10}; it is not a separate existence technique.

\begin{theorem}[Supporting lower bound]\label{thm:lower}
For every $k\geq1$, a finite $(k+2)$-player reachability game with $O(k^2)$
vertices, outdegree at most two, and no cycles except terminal loops admits
an SPE winning for a designated player $E$, but every such SPE requires at
least $\binom{k}{\lfloor k/2\rfloor}\geq2^k/(k+1)$ vertex-update states for
$E$, irrespective of the other players' memory.
\end{theorem}

Consequently the worst-case memory for constrained pure SPEs is
$2^{\Theta(n)}$. Exponential SPE memory necessity and its checking mechanisms
already have antecedents~\cite{BGMR23,FH12}. The lower bound is included to
make the quantitative limitation self-contained; no priority claim for
exponential necessity or the precise lower-bound constants is made.

\subsection{Prior work and the precise advance}

The finite-arena results of Brihaye, Bruy\`ere, Goeminne and
Raskin~\cite{BBGR18,BBGR21} provide the closest reconstruction baseline.
In their full version~\cite{BBGR18}, Proposition~9 identifies weak
SPEs with SPEs for reachability, and Proposition~19 and Corollary~21 give
arena-dependent payoff-preserving reconstruction through finite symbolic
witnesses. Their Theorem~34 already establishes fixed-parameter tractability
of constrained existence. We claim no new FPT classification.

A close antecedent is Bruy\`ere, Le Roux, Pauly and
Raskin~\cite{BLPR21}. Theorem~7 of their author postprint uses finitely many
leaf outcomes and outcome-indexed positional paths to construct finite-memory
weak SPEs under reachable absorbing-leaf hypotheses. Theorem~16 gives
finite-range weak-SPE existence. General reachability arenas do not satisfy
the absorbing-leaf hypotheses: completing all targets in a promised payoff
does not eliminate legal future deviations by the remaining players.
We credit outcome memory and positional path selection to this antecedent.

Main~\cite{Main26}, Theorem~23 of the author version, gives quadratic
player-dependent memory bounds for Nash equilibria in reachability games.
Its Remark~1 distinguishes vertex from edge memory. Nash stability alone
does not provide credible continuations after every history. Flesch,
Kuipers, Mashiah-Yaakovi, Schoenmakers, Solan and Vrieze~\cite{FKMSSV10},
Section~3.3, already use payoff-minimizing continuations after deviations.
Alluwaym, Main and Schewe~\cite{AMS26} give stationary randomized SPEs on
finite graphs and a pure positional special case for absorbing targets;
this does not supply prescribed-payoff pure compression on arbitrary arenas.

The advance here is a compatibility and previous-owner normalization rule
for all-history SPE payoff labels on continuing arenas. A fixed label is
realized by a progress table depending on the current vertex and visited
set. A promise is retained whenever it remains available. If it disappears,
the previous owner's payoff bit determines a credible reset. Storing that
owner requires $n$ states rather than a register ranging over arena vertices.
An invariant for every unilateral deviation proves subgame perfection.
The inspected primary literature supports this precise distinction;
the comparison is bounded and is not an exhaustive priority certificate.

\section{Model and memory convention}\label{sec:model}

Let $N=\{1,\ldots,n\}$ and let $V$ be a set. An arena consists of a relation
$\mathcal E\subseteq V\times V$, an owner map $\own:V\to N$, and an initial
vertex $v_0$. Every $v$ has a nonempty set $\Succ(v)$ of successors.
Player $i$ has a target $T_i\subseteq V$. Put $t(v)=\{i:v\in T_i\}$.
A legal history is a finite nonempty path from $v_0$; the set of these
histories is $\Hist$. A play is an infinite path. Its payoff set consists
of the players whose targets occur at some finite position. We identify a
payoff set $p\subseteq N$ with its Boolean vector, writing $p_i=1$ iff $i\in p$.

A pure strategy chooses a successor after each history ending at a vertex
owned by its player. A profile $\sigma$ fixes a unique continuation after
each $h\in\Hist$. Write $p_\sigma(h)$ for the payoff set of the whole
history followed by this continuation, counting their boundary vertex once.
In particular target visits before a subgame remain satisfied. A profile is
an SPE if, after every $h\in\Hist$, no player can improve its Boolean payoff
by any unilateral strategy change. Deviating strategies can use arbitrary
memory and change their moves infinitely often.

The memory convention is
\[
 M,\quad m_0\in M,\quad \up:M\times V\to M,
 \quad \nxt_i:M\times V_i\to V,
\]
where $V_i=\own^{-1}(i)$ and $\nxt_i(m,v)\in\Succ(v)$. At a history $hv$,
where $h$ may be empty, the strategy plays
\[
                    \nxt_i(\up^*(m_0,h),v).
\]
Thus the incoming memory has processed the prefix strictly before $v$.
The next-move and update functions both receive this incoming state and
the current vertex. No previous vertex is available unless explicitly
stored. On an infinite arena these functions may be arbitrary tables;
only $|M|$ is counted. We work with ordinary set-theoretic choice.

\section{Credible labels and uniform progress}\label{sec:labels}

Fix the SPE $\sigma$ to be compressed. For a history $h$, put
$W(h)=\bigcup_{v\text{ in }h}t(v)$ and $x(h)=(\last(h),W(h))$.
Let $X=\{x(h):h\in\Hist\}$. An augmented state $x=(v,W)$ has, for every
$w\in\Succ(v)$, augmented successor $y=(w,W\cup t(w))\in X$.
Define
\[
 P(x)=\{p_\sigma(h):h\in\Hist,\ x(h)=x\}.
\]
The label set $P(v,W)$ is a nonempty finite subset of $2^N$, and each
of its labels contains $W$. It aggregates all legal histories, not just
the histories consistent with $\sigma$. Neither a witnessing history nor
the original strategy's memory will be stored.

\begin{lemma}[Local credibility]\label{lem:local}
For $x=(v,W)\in X$, $p\in P(x)$, $i=\own(v)$, and every augmented
successor $y$ of $x$, some $q\in P(y)$ satisfies $q_i\leq p_i$.
\end{lemma}
\begin{proof}
Choose a history $h$ witnessing $p\in P(x)$ and let $w$ be the vertex
component of $y$. The original continuation after $hw$ gives a label
$q=p_\sigma(hw)\in P(y)$. Moving to $w$ at $h$ and thereafter following
$\sigma_i$ is an admissible unilateral change after $h$. Since $\sigma$
is subgame perfect, its payoff $q_i$ cannot exceed $p_i$.
This also covers the prescribed successor, for which $q=p$.
\end{proof}

\begin{lemma}[Progress without a rank register]\label{lem:progress}
For every $p\subseteq N$, there is a successor table $f_p$ on
$X_p=\{x\in X:p\in P(x)\}$ such that, from every $x=(v,W)\in X_p$,
following $f_p$ stays in $X_p$ and realizes exactly payoff $p$, including $W$.
\end{lemma}
\begin{proof}
Choose a history witnessing $p$ at $x$ and follow $\sigma$. Every augmented
state on that continuation is in $X_p$. Each player in $p\setminus W$
visits its target after finitely many edges. Since there are finitely many
players, the maximum of these finitely many visit times is finite (use zero
if $p\setminus W$ is empty).
Thus some finite path in $X_p$ reaches
\[
                  C_p=\{(u,B)\in X_p:B=p\}.
\]
For $x\in X_p$, let $d_p(x)$ be the minimum length of such a path.
It is a finite natural number, with no uniform bound asserted. If
$d_p(x)>0$, choose $f_p(x)$ to give a successor $y\in X_p$ with
$d_p(y)=d_p(x)-1$. If $d_p(x)=0$, choose any successor in $X_p$, which
the witnessing continuation supplies. At rank zero $W=p$, and a successor
in $X_p$ has visited set containing $W$ and contained in $p$, hence equal
to $p$. The selected path therefore reaches rank zero in finitely many
steps and thereafter visits no target outside $p$. Its full payoff is $p$.
The choices are fixed in the vertex-dependent table, not in a memory
register. Taking a minimum over natural numbers requires no finite branching.
\end{proof}

\section{The vertex-update construction}\label{sec:machine}

Let $p_0=p_\sigma(v_0)$. Use the common memory set
\[
 M=\{(A,p,j): A\subseteq p\subseteq N,\ j\in N\},
 \qquad m_0=(\varnothing,p_0,1).
\]
For each player there are three possibilities in the pair $(A,p)$:
outside $p$, inside $p\setminus A$, or inside $A$. Thus $|M|=n3^n$.
The set $A$ records visits strictly before the current vertex, $p$ is an
incoming promise, and $j$ is the owner of the previous processed vertex.

Fix a total order on $2^N$. For $x\in X$ let $b_j(x)$ be the first label
in $P(x)$ minimizing its $j$th bit. Given $m=(A,p,j)$ and a vertex $v$,
form $B=A\cup t(v)$ and $x=(v,B)$. If $x\in X$, define
\[
 R(m,v)=\begin{cases}
 p,&p\in P(x),\\
 b_j(x),&p\notin P(x).
 \end{cases}
\]
Writing $r=R(m,v)$, set
\begin{equation}\label{eq:machine}
 \up(m,v)=(B,r,\own(v)),\qquad
 \nxt_i(m,v)=f_r(x)\quad(v\in V_i).
\end{equation}
Both functions use the same incoming $m$: the update does not process a
successor, and the next-move function does not receive already updated
memory. Since $r\in P(x)$, we have $B\subseteq r$, so the updated state
belongs to $M$. If $x\notin X$, complete both functions arbitrarily with
a state in $M$ and a legal successor. These arguments are unreachable from
the initialized machine on a legal history. Call the profile $\tau$.
Every player uses the same update, so their memories agree on
every history, even when some players have deviated.

\begin{lemma}[All-history realization]\label{lem:realization}
At every legal history $hv$, the incoming $A$ is precisely the visited set
of $h$, and the normalized promise $r=R(m,v)$ belongs to $P(v,A\cup t(v))$.
Following $\tau$ after $hv$ gives full payoff exactly $r$.
\end{lemma}
\begin{proof}
Initially $A=\varnothing$. The update in~\eqref{eq:machine} stores exactly
$A\cup t(v)$, independently of the successor taken. Induction over every
legal history gives the first assertion and ensures $x\in X$ at each
normalization. By definition the normalized label is in $P(x)$.

If the prescribed move $f_r(x)$ is taken, its augmented successor $y$
belongs to $X_r$ by Lemma~\ref{lem:progress}. The incoming promise there
is $r\in P(y)$, so normalization retains it. Continuing to follow $\tau$
retains $r$ forever and follows $f_r$. Lemma~\ref{lem:progress} gives full
payoff $r$, including all visits in the prefix.
\end{proof}

\begin{proof}[Proof of Theorem~\ref{thm:main}]
Fix any legal history $H$ and player $i$, and let $r$ be the current
normalized promise. By Lemma~\ref{lem:realization}, $r_i$ is the actual
payoff under $\tau$ after $H$. If $r_i=1$, no Boolean improvement is possible.
Suppose $r_i=0$ and allow any unilateral deviating strategy for $i$.
We show that every normalized promise along its continuation still excludes $i$.

At another player's vertex the prescribed move is taken, retaining the
promise at the next vertex. At a vertex owned by $i$, let $y$ be the
augmented state of the successor chosen by the deviation. The update at
the current vertex stores previous owner $i$. If the old promise is in
$P(y)$, it is retained. Otherwise the new promise is $b_i(y)$.
Lemma~\ref{lem:local} supplies some $q\in P(y)$ with $q_i\leq r_i=0$,
so the minimizing label also has bit zero. This proves the invariant by
induction at every finite position, even with infinitely many deviations.

Every normalized label contains every target visited so far. Hence $T_i$
is never visited on the deviating play, and its payoff is zero. No player
can profitably deviate after $H$. As $H$ was any legal history, $\tau$ is
an SPE. Finally $p_0\in P(v_0,t(v_0))$, so the initial normalization retains
$p_0$ and Lemma~\ref{lem:realization} proves exact initial payoff preservation.
The machine has $n3^n$ states, as required.
\end{proof}

\section{Unconditional existence from an earlier theorem}\label{sec:existence}

\begin{proof}[Proof of Corollary~\ref{cor:existence}]
Theorem~2.3 of~\cite{FKMSSV10} gives a pure subgame-perfect
$\varepsilon$-equilibrium for every $\varepsilon>0$ in finite-player
perfect-information games with bounded lower-semicontinuous payoffs.
Its model allows arbitrary action sets (Definition~2.1 and footnote~2),
history-dependent feasible actions (Remark~2.1), and pure strategies without
measurability constraints (Definition~2.2). Its subgame payoffs retain the
prefix (Definition~2.4).

Unfold our arena into its legal-history tree. At a history ending in $v$,
the active player is $\own(v)$ and the available actions are $\Succ(v)$.
This is a nonempty set, without a cardinality restriction. One may use $V$
as the common action alphabet and restrict feasible actions by history.
The Boolean payoff for visiting $T_i$ is bounded and lower semicontinuous:
its one-set is open in the topology generated by finite-history cylinders,
because a finite target visit witnesses victory. This holds also for an
uncountable discrete alphabet and includes initial target visits.
Apply the old theorem with $\varepsilon=1/2$. A profitable Boolean deviation
would improve payoff by one, contradicting its inequality in that subgame.
Thus the obtained profile is an exact pure SPE. Theorem~\ref{thm:main}
compresses it with the claimed memory bound.
\end{proof}

The existence input is essential to this particular deduction and is fully
credited to~\cite{FKMSSV10}. The compression proof itself begins with any
given SPE and does not require an algorithm for finding one or computing $P(x)$.

\section{A constrained exponential lower bound}\label{sec:lower}

We give the complete family for Theorem~\ref{thm:lower}. Its ingredients
are related to generalized-reachability subset memory~\cite{FH12} and to
the credible checking constructions surveyed in~\cite[Section~3.1.3]{BGMR23}.
The argument applies to every SPE satisfying a required winner, not just
to profiles carrying a particular prescribed payoff.

Fix $k\geq1$ and put $r=\lfloor k/2\rfloor$. The players are
$A,E,C_1,\ldots,C_k$. First $A$ selects exactly $r$ indices from $[k]$,
giving a set $W$. Then $E$ selects exactly $k-r$ indices, giving a set $U$.
Each phase examines indices in order and merges all its histories at its
end. In particular every $r$-subset $W$ reaches the same first $E$ vertex.

Here is an explicit implementation. A quota-$b$ phase has decision vertices
$q(i,c)$, with $0\leq i<k$, $0\leq c\leq i$ and $c\leq b\leq c+k-i$.
The two phases use disjoint copies. At $q(i,c)$ the selecting player may
choose $d\in\{0,1\}$ exactly when
\[
                         c+d\leq b\leq c+d+k-i-1.
\]
The edge goes to a forced marker $z(i,c,d)$ and then to $q(i+1,c+d)$,
or to the next phase when $i+1=k$. A marker with $d=1$ is an inclusion
of index $i+1$. Every decision has one or two legal choices. The count is
an arena vertex component, not a memory state. There are $O(k^2)$ vertices.

After both phases, play visits checking vertices $c_1,\ldots,c_k$.
At $c_i$, player $C_i$ may go to terminal $a$ or continue to $c_{i+1}$
(to terminal $e$ when $i=k$). Both terminals have only self-loops.
The targets of $A$ and $E$ are $\{a\}$ and $\{e\}$, respectively.
The target of $C_i$ consists of $a$ and all inclusion markers of index $i$
in either phase. There are no other target memberships. The initial
vertex is the first $A$ decision. Ownership at forced vertices is immaterial.
The arena is finite, has no dead ends, and has outdegree at most two.

\begin{lemma}[Credible terminal outcome]\label{lem:checker}
In any pure SPE, after any legal history reaching the checking phase with
$W\cup U\ne[k]$, the continuation ends at $a$.
\end{lemma}
\begin{proof}
Otherwise it ends at $e$, passing every checking vertex. A checker $C_i$
with $i\notin W\cup U$ then receives zero, but can move from $c_i$ to $a$
and receive one. This contradicts subgame perfection at that history.
The finite checking chain has no other outcome.
\end{proof}

This lemma does not require the first unsatisfied checker to exit
immediately; a later exit could also give it payoff one.

\begin{proof}[Proof of Theorem~\ref{thm:lower}]
An $E$-winning SPE exists as follows. Player $A$ uses any fixed legal
selection. Player $E$ remembers $W$ and, at index $i$, tries to include
$i$ exactly when $i\notin W$; if previous off-profile choices make this
choice illegal, it takes any legal choice. At $c_i$, player $C_i$ continues
if its target has already been visited and otherwise goes to $a$.

After any history in $A$'s phase, all its possible completed choices produce
an $r$-subset, to which $E$ responds with its complement. All checkers
then win and continue to $e$. Player $A$ loses regardless of any unilateral
change; all other players win. After a history in $E$'s phase, if its past
choices agree with the complement, the unique complementary completion
is legal and leads to $e$. If a past choice disagreed, covering all indices
is impossible: an omitted complement index cannot be revisited, while an
included index of $W$ consumes one of the exactly $k-r$ response slots.
A missing index then causes a checker to exit, so $E$ cannot improve from
that subgame. Player $A$ has no further choices. In the checking phase,
a satisfied checker already wins, and an unsatisfied checker wins by exiting.
A checker whose decision is past cannot affect its continuation. These
observations cover arbitrary earlier checking deviations and terminal
histories too. Thus the profile is an SPE after every legal history.
Its initial outcome is $e$. It has finite vertex-update memory: $E$ stores
$W$, each checker stores its own visited-target bit, and $A$ is positional.

Now fix any SPE winning initially for $E$, with no memory restrictions
on any player. Its initial outcome is $e$, so $A$ gets zero. Player $A$
can unilaterally select any $r$-subset $W$. Holding all other strategies
fixed, the resulting outcome must still give $A$ zero. If the prescribed
response $U$ failed to cover $[k]$, Lemma~\ref{lem:checker}, applied to
that legal off-profile history, would force $a$ and yield a profitable
initial deviation for $A$. Therefore the single fixed strategy of $E$
must respond to every $W$ by $U=[k]\setminus W$.

Suppose a vertex-update machine represents this strategy of $E$. Let $m(W)$
be its incoming memory at the common first $E$ vertex after $A$ selects $W$.
If $m(W)=m(W')$ for distinct $W,W'$, determinism forces the same response
vertex sequence: the current vertex and incoming state agree, all other
moves during this phase are forced, and the update sees the same successive
vertices. Hence the response sets coincide. They cannot equal both distinct
complements. Thus $W\mapsto m(W)$ is injective, giving at least
$\binom{k}{r}$ states. Other players cannot signal through their memories
during the response phase. The central binomial coefficient is the largest
of $k+1$ coefficients summing to $2^k$, giving the stated inequality.
The case $k=1$ is included; its first quota is zero and its bound is one.
\end{proof}

\section{Verification and limits}

The proofs establish the arbitrary-arena upper bound and the all-$k$
lower bound; finite computations are supplementary checks. Exact Python
best-response checks were applied after every legal-history product state,
allowing all edges to generate histories and arbitrary unilateral
deviations against fixed opponents. A separate implementation checked
2,827 verified finite SPE inputs and their recompressions, preserving exact
initial payoffs; it also checked the lower family for $k=1,\ldots,6$.
Deliberately omitting the previous-owner register or the progress condition
produced failing controls. These tests are neither an exhaustive enumeration
of arbitrary games nor formal proof-assistant certification.

Empty targets, initially visited targets, and empty or full payoff sets
are included. Finitely many players is used to fulfill a promise after
finitely many target visits. No bound on the selector's distance is needed.
The result supplies no effective algorithm for the arbitrary infinite
tables, no new finite-arena constrained-existence FPT result, and no extension
to quantitative arrival-time, B\"uchi, or arbitrary omega-regular objectives.
The gap between the exponential bases of the upper and lower bounds remains.

\section*{AI methodology}
OpenAI Codex was used substantially in the research and preparation of this
article, including proof construction and adversarial checking, the explicit
memory machine, the existence-theorem application, the lower-bound family,
literature comparison, exact Python verification, and LaTeX drafting.
These uses materially affected all technical sections. Additional fresh
model contexts checked the mathematical arguments and primary-source
comparisons and produced a separate finite checker. Such model checks are
not external peer review or formal certification. This disclosure does not
assert that the listed author has already performed a separate human
verification of every generated argument or reference. Human verification
and responsibility are required before any external submission.

\clearpage
\bibliographystyle{alpha}
\bibliography{references}
\end{document}
